\begin{textsample}{POS Dim 1 – ms – Score 26.00 – ms/ms\_em\_30.txt}  \label{ex:f1_pos_208}
3.1.1.
Seriação da Área de Linguagens e suas Tecnologias As habilidades da área de Linguagens e suas Tecnologias estão organizadas nos três anos do Ensino Médio, considerando o grau de complexidade e a carga horária necessária ao seu desenvolvimento. 3.2.
Área de Conhecimento : Matemática e suas Tecnologias A área de conhecimento Matemática e suas Tecnologias, assim definida pela LDB, é desenvolvida na etapa do Ensino Médio com foco em competências e habilidades direcionadas para formação geral e contextualizada do estudante, conforme artigo 35 : > O ensino médio, etapa final da educação básica, com duração mínima de três anos, terá como finalidades : > I - a consolidação e o aprofundamento dos conhecimentos adquiridos no ensino fundamental, possibilitando o prosseguimento de estudos. > II - a preparação básica para o trabalho e a cidadania do \textbf{educando}, para continuar aprendendo, de modo a ser capaz de se adaptar com flexibilidade a novas condições de ocupação ou aperfeiçoamento posteriores. > III - o aprimoramento do \textbf{educando} como pessoa humana, incluindo a formação ética e o desenvolvimento da autonomia intelectual e do pensamento crítico ( BRASIL, 1996 ).
Assim, o currículo da área de Matemática e suas Tecnologias também objetiva a ampliação do letramento matemático ao propor metodologias de aprendizagem \textbf{que} contribuam com a formação global do estudante, \textbf{uma} vez \textbf{que} o coloca frente a situações reais envolvendo objetos de conhecimento e tecnologias da informação para posicionar -se de forma crítica e exercer o protagonismo ao compreender e solucionar problemas da realidade.
Conforme citado pelo Programa Internacional de Avaliação de Estudantes ( Pisa ) no Brasil : > O letramento matemático refere -se à capacidade de identificar e compreender o papel da Matemática no mundo \textbf{moderno}, de tal forma a fazer julgamentos bem \textbf{embasados} e a utilizar e envolver -se com a Matemática, com o objetivo de atender às necessidades do indivíduo no cumprimento de seu papel de cidadão consciente, crítico e construtivo ( BRASIL, 2010b ).
Dessa forma, propõe -se \textbf{que} os estudantes construam \textbf{uma} visão mais integrada da Matemática em relação aos conhecimentos já explorados e compreendam sua importância para a sociedade, abrangendo também os temas \textbf{contemporâneos} de maneira transdisciplinar.
Na Base Nacional Comum Curricular ( BNCC ), a área de Matemática e suas Tecnologias está organizada, assim como no Ensino Fundamental, em eixos temáticos : Números e Álgebra, Geometria e Medidas, Probabilidade e Estatística.
Porém, a aprendizagem da Matemática no Ensino Médio deve estimular processos mais elaborados de reflexão e de abstração e permitir ao estudante a investigação, construção de modelos, formulação e resolução de problemas em diversos contextos ( do cotidiano, da própria Matemática e de outras áreas do conhecimento ) de forma autônoma.
A BNCC cita as tecnologias digitais como prática de aprendizagem e permite acrescentar novas possibilidades, \textbf{que} favoreçam o processo de desenvolvimento do estudante, na sua autonomia, forma de agir e transformar a sociedade, além de desenvolver suas potencialidades e estimular a criatividade.
Segundo D’Ambrosio : > Ao longo da evolução da \textbf{humanidade}, Matemática e tecnologia se desenvolveram em \textbf{íntima} associação, numa relação \textbf{que} poderíamos \textbf{dizer} simbiótica.
A tecnologia entendida como convergência do saber ( ciência ) e do fazer ( técnica ), e a matemática são intrínsecas à busca solidária do sobreviver e de transcender.
A geração do conhecimento matemático não pode, portanto ser dissociada da tecnologia disponível ( D’AMBROSIO 1996, p. 13 ).
Dessa forma, o uso das tecnologias digitais da informação deve estar ligado à aprendizagem da Matemática, pois grande parte dos estudantes está inserida no mundo digital, é ativa em redes sociais, produz vídeos, podcast, vlogs, dentre outros. \textbf{Logo}, o \textbf{educando} deve ser também protagonista no estudo da Matemática com o uso da internet ou de softwares ( Apps ).
Assim, é possível \textbf{uma} investigação matemática, juntamente com o pensamento computacional, por meio da interpretação e elaboração de algoritmos, o \textbf{que} permite aos estudantes \textbf{um} aprofundamento e \textbf{uma} ampliação de sua aprendizagem, com a consolidação de sua autoestima e autonomia.
Anteriormente, os conteúdos listados nos currículos eram os objetos da prática docente, no entanto, na BNCC, são apresentadas cinco competências específicas a serem desenvolvidas, cada \textbf{uma} com sua lista de habilidades, \textbf{que} representam aprendizagens a serem garantidas a cada estudante.
Para tal progresso, na área de Matemática e suas Tecnologias, devem -se provocar ações \textbf{que} façam o estudante observar, modificar e evoluir o modo próprio de raciocinar, representar, comunicar e \textbf{argumentar}.
Fato \textbf{que} se dá pela interação do estudante com seus colegas e professores, a partir da elaboração de registros de representações \textbf{que} possam impulsionar a modelagem de situações diversas por meio da linguagem específica da Matemática na busca de soluções e respostas a problemas, formulações e comprovações de conjecturas, dentre outras situações, mediante o uso de argumentação consistente na justificativa de resultados.
A BNCC propõe \textbf{que} a aprendizagem da Matemática passe a ter \textbf{uma} aplicação na vida real, com \textbf{estímulo} à produção/aquisição de novos conhecimentos, de forma significativa, promovendo o protagonismo juvenil dos estudantes.
Além disso, \textbf{que} essa aprendizagem privilegie o desenvolvimento da dimensão afetiva e cognitiva com ações \textbf{que} colaborem para a formação da cidadania, ou seja, preparar esses estudantes para os desafios da sociedade \textbf{contemporânea}.
Nesse \textbf{raciocínio}, o processo de aprendizagem deve buscar o desenvolvimento de todas as dimensões do ser humano.
A área de Matemática e Suas Tecnologias é composta por quarenta e cinco habilidades ligadas a competências específicas representadas por código alfanumérico, mas isso não significa \textbf{uma} dissociação das habilidades de competências diferentes, nem \textbf{que} \textbf{uma} habilidade de determinada competência não contribua para o desenvolvimento de outra.
Conforme a BNCC, “ A área de Matemática e suas Tecnologias deve garantir aos estudantes o desenvolvimento de competências específicas.
Relacionadas a cada \textbf{uma} delas, são indicadas, posteriormente, habilidades a ser alcançadas nessa etapa ” ( BRASIL, 2018c, p. 528 ).
As competências e as habilidades na área de Matemática e suas Tecnologias formam \textbf{um} todo conectado, de modo \textbf{que} o desenvolvimento de \textbf{uma} requer, em determinadas situações, a mobilização de outras.
Cabe observar \textbf{que} essas competências consideram \textbf{que}, além da \textbf{cognição}, os estudantes devem desenvolver atitudes de autoestima, de perseverança na busca de soluções e de respeito ao trabalho e às opiniões dos colegas, mantendo predisposição para realizar ações em grupo ( BRASIL, 2018c. p. 530 ).
A primeira competência utiliza a Matemática para interpretar situações em diversos contextos e fazer julgamentos bem fundamentados, prevê \textbf{uma} formação científica geral e capacidade de analisar criticamente o \textbf{que} é divulgado no meio de comunicação.
A competência dois propõe aos estudantes investigações de caráter social, de modo a propor ou participar na solução de problemas, com base em princípios solidários, éticos e sustentáveis, refletindo sobre a importância da Matemática no contexto sociopolítico e cultural.
A terceira competência \textbf{salienta} a utilização da Matemática para interpretar, construir modelos e resolver problemas com a análise da plausibilidade dos resultados e adequações das soluções propostas, de modo a construir argumentação consistente.
Nessa competência, o estudante deverá ser conduzido a pensar na resolução e formulação de problemas ; conforme a BNCC enfatiza, deve -se substituir “ Resolver Problemas ” por “ Resolver e Elaborar Problemas ”.
Na competência quatro, o estudante deve compreender e utilizar diferentes registros de representação matemática na busca de solução e comunicação de resultados de problemas.
Nesse caso, o estudante deve ser estimulado a explorar mais de \textbf{um} registro de representação, sempre \textbf{que} possível.
Ao representar matematicamente diferentes formas de \textbf{um} mesmo resultado, o estudante potencializará diferentes formas de resolver problemas, intensificando a sua capacidade matemática.
A quinta e última competência, \textbf{que} é investigar e estabelecer conjecturas a respeito de diferentes conceitos e propriedades matemáticas, empregando estratégias e recursos, como observação de padrões, experimentações e diferentes tecnologias, identificando a necessidade, ou não, de \textbf{uma} demonstração cada vez mais formal na validação das referidas conjecturas, estabelece \textbf{que} o estudante deve formular hipóteses e teses com base em suas investigações, buscar contraexemplos e, quando necessário, \textbf{argumentar} a fim de prová -las.
É \textbf{imprescindível} \textbf{que} o \textbf{sujeito} conviva com o \textbf{raciocínio} hipotético-dedutivo, diferente das outras ciências \textbf{que} utilizam o \textbf{raciocínio} hipotético-indutivo.
Assim, é possível observar certa progressão nas cinco competências específicas da área de Matemática e suas Tecnologias.
As habilidades relacionadas em cada competência preveem objetos de conhecimento \textbf{que} contribuem para o estudante adquirir e aperfeiçoar determinada competência, tanto no âmbito cognitivo quanto no socioemocional.
Verbos, tais como : interpretar, analisar, identificar, propor, participar, aplicar, resolver, elaborar, empregar, investigar, representar e reconhecer são comuns nas habilidades matemáticas.
Essas habilidades não representam, de maneira \textbf{explícita}, os conteúdos de matrizes, binômio de Newton, números complexos, polinômios e geometria analítica, porém nada impede a \textbf{abordagem} desses assuntos, \textbf{dependendo} do tema \textbf{que} determinada pesquisa ou estudo \textbf{necessitar}.
Em contrapartida, outros temas ganham maior visibilidade, tais como : ladrilhamento, educação financeira, planilhas eletrônicas, fluxograma, linguagem de programação, interpretação de texto científico, análise de taxas, planejamento e execução de pesquisas, ações envolvendo a utilização de aplicativos, tecnologias matemáticas digitais e projeções \textbf{cartográficas}, \textbf{que} são temas \textbf{atuais} e fundamentais, os quais \textbf{necessitam}, algumas vezes, de pesquisas e atualizações tanto para o estudante como para o professor.
Destaca -se \textbf{que} o quadro organizador contém as cinco competências específicas previstas na BNCC, cada \textbf{uma} com seu conjunto de habilidades, eixos temáticos, objetos de conhecimentos e sugestões didáticas.
É possível mesclar as habilidades de cada competência específica e a ordem em \textbf{que} elas aparecem.
Assim, o modo e a sequência em \textbf{que} se aborda cada habilidade \textbf{depende} da realidade de cada unidade escolar.
As sugestões didáticas descrevem possíveis ações ou condutas esperadas na prática docente, visando \textbf{fomentar} o protagonismo e a autoria do estudante ; o professor não deve tomá -las como se fossem \textbf{um} currículo a seguir, pois essas sugestões apenas norteiam a elaboração de sua metodologia.
Ressalta -se \textbf{que} certas sugestões didáticas são mais adequadas a determinadas regiões do Estado, cabendo ao docente fazer as adaptações necessárias.
É importante \textbf{notar} também \textbf{que} a sugestão didática nem sempre contempla todos os objetos de conhecimento, \textbf{necessitando}, assim, de outras ações.

% matched lemmas: abordagem, argumentar, atual, cartográfico, cognição, contemporâneo, depender, dizer, educar, embasar, estímulo, explícito, fomentar, humanidade, imprescindível, logo, moderno, necessitar, notar, que, raciocínio, salientar, sujeito, um, uma, íntimo
\end{textsample}
